\documentclass[11pt]{article}
\usepackage[margin=1in]{geometry}
\usepackage{amsmath,amssymb,amsthm}
\usepackage{booktabs}
\usepackage[hidelinks]{hyperref}

\theoremstyle{plain}
\newtheorem{theorem}{Theorem}
\newtheorem{proposition}[theorem]{Proposition}
\newtheorem{lemma}[theorem]{Lemma}
\newtheorem{corollary}[theorem]{Corollary}
\theoremstyle{definition}
\newtheorem{definition}[theorem]{Definition}
\theoremstyle{remark}
\newtheorem{remark}[theorem]{Remark}

\newcommand{\N}{N}
\newcommand{\g}{\mathfrak{g}}

\title{\bf Cute Goldbach Gaps}
\author{Anonymous}
\date{\today}

\begin{document}
\maketitle

\begin{abstract}
For an even integer $\N$, let $\g(\N)$ be the least prime appearing in a
Goldbach partition of $\N$---the quantity $p(n)$ of the \emph{minimal Goldbach
partition} studied by Granville, van de Lune and te Riele, who conjectured
$p(n)=O((\log n)^2\log\log n)$ and for which $p(n)\le 9781$ holds empirically for
all even $n\le4\cdot10^{18}$. We ask how large $\g(\N)$ can be \emph{forced}, and
at what cost in the size of $\N$. After an elementary warm-up showing that the
neighbourhood of a factorial inflates $\g$ via divisibility and Wilson's theorem
(for $\N=27!+28$ we get $\g(\N)=71$), we prove that $\g$ is \emph{unbounded}, and
we do so unconditionally: combining a primorial obstruction with Dirichlet's
theorem produces, for every prime $P$, an even $\N$ with an \emph{explicit}
Goldbach partition and $\g(\N)=\mathrm{nextprime}(P)$. We then show that a
covering system of congruences forces $\g(\N)$ above any target $Q$ using a
modulus exponentially smaller than the primorial, and that any finite covering
can likewise be realised unconditionally by a double application of Dirichlet.
We exhibit a $1020$-digit even integer whose least Goldbach summand is
$\g(\N)=100\,747$, its prime complement proved prime by an Atkin--Morain
elliptic-curve certificate, and we certify a $327\,455$-digit even integer for
which no prime below $10^8$ is a summand. To our knowledge these are,
respectively, the largest \emph{proven} least Goldbach summand on record (about
ten times the previous record of $9781$ near $3.3\cdot10^{18}$) and the largest
explicit Goldbach-summand obstruction yet reported.
Finally we record the
\emph{one-directional} transfer from large prime gaps: a long prime gap forces a
large least summand, but the converse fails, since only prime offsets need to be
covered---a sparser problem. The verified integer, the congruence system that
defines it, and a standalone verification script with SHA-256 hashes are given in
the appendices.
\end{abstract}

\section{Introduction}

Goldbach's conjecture asserts that every even integer $\N>2$ is a sum of two
primes; it has been verified for all $\N\le4\cdot10^{18}$~\cite{OeS}.
Heuristically one expects not just one representation but many. The
Hardy--Littlewood circle method---which estimates such counts by treating the
primes as a random set of the appropriate density and then correcting for
divisibility constraints---predicts that the number $r(\N)$ of ways to write
$\N=p+q$ with $p,q$ prime is
\begin{equation}\label{eq:HL}
r(\N)\;\sim\;2\,\Pi_2\,\frac{\N}{(\log \N)^2}\prod_{\substack{p\mid \N\\ p>2}}\frac{p-1}{p-2},
\qquad
\Pi_2=\prod_{p>2}\Bigl(1-\tfrac{1}{(p-1)^2}\Bigr).
\end{equation}
Reading this from left to right: $\N/(\log \N)^2$ is the naive guess, since a
number near $\N$ is prime with probability about $1/\log \N$ (the prime number
theorem) and there are $\N$ choices for $p$; the constant
$\Pi_2=\prod_{p>2}(1-\tfrac1{(p-1)^2})\approx0.6601$---an infinite product running
over all odd primes $p$---is the standard correction (the twin-prime constant)
that accounts for primes not being truly random mod small numbers; and the last
factor $\prod_{p\mid \N,\,p>2}\frac{p-1}{p-2}$, a product over the odd primes
\emph{dividing} $\N$, inflates the count when $\N$ is highly composite. The point
for us is only the conclusion: $r(\N)\to\infty$, so large even integers have
\emph{many} Goldbach partitions, and the interesting question is not whether a
representation exists but how lopsided the most lopsided one must be.

\begin{definition}
For an even integer $\N\ge4$ admitting a Goldbach representation, the
\emph{least Goldbach summand} is
\[
\g(\N)\;=\;\min\{\,p\ \text{prime}:\N-p\ \text{is prime}\,\}.
\]
\end{definition}

This is exactly the quantity $p(n)$ attached to the \emph{minimal Goldbach
partition} $\N=p+q$ (with $p$ minimal) by Granville, van de Lune and te
Riele~\cite{GLvdtR}, who conjectured that it never grows faster than a fixed
multiple of $(\log \N)^2\log\log \N$:
\begin{equation}\label{eq:GLR}
\g(\N)\;=\;O\bigl((\log \N)^2\log\log \N\bigr).
\end{equation}
The notation $f=O(g)$ means $|f|\le C\,g$ for some constant $C$ and all large
$\N$; here \eqref{eq:GLR} is a \emph{conjectural} ceiling, not a theorem, and one
of the things we do below is locate it between provable bounds.
Empirically $\g$ is tiny and grows very slowly: its record values are catalogued
in the OEIS~\cite{OEIS} (sequence A020481, with records A025018 and A025019).
Exhaustive verification has produced a slowly climbing ladder of records---among
them $\g=1093$ at $\N=60{,}119{,}912$, then $\g=5569$ for all even
$\N\le4\cdot10^{14}$~\cite{Richstein}, and the current record $\g=9781$, first
attained at $\N=3{,}325{,}581{,}707{,}333{,}960{,}528$, with $\g(\N)\le9781$ for
\emph{all} even $\N\le4\cdot10^{18}$~\cite{OeS}. (We reconfirmed the first and
last of these deterministically, both being below $2^{64}$.) Strikingly, these
records occur at integers of at most $19$ digits. Our question
is the opposite, extremal and constructive:
\begin{quote}\itshape
How large can $\g(\N)$ be forced to be, and at what cost in the size of $\N$?
\end{quote}
A companion quantity is $S(q)$, the least even $\N$ with $\g(\N)=q$~\cite{GLvdtR}:
the ``smallest Goldbach desert'' for a prescribed least summand.

\paragraph{Results and roadmap.} We give three constructions of increasing power,
and are careful throughout to separate three distinct claims: \emph{(a)} no small
prime is a summand of $\N$; \emph{(b)} $\N$ has a Goldbach representation at all;
and \emph{(c)} the least summand equals a stated value. A covering argument gives
\emph{(a)} unconditionally; \emph{(b)} requires either Goldbach's conjecture, an
exhibited prime complement, or a separate existence argument; and \emph{(c)}
requires both.
\begin{itemize}
\item[\S\ref{sec:factorial}] \textbf{Factorial neighbourhoods.} Near $n!$,
divisibility kills $\N-q$ for small $q$, and Wilson's theorem adds one
obstruction; for $\N=27!+28$ every prime $q\le67$ fails, so $\g(\N)=71$.
\item[\S\ref{sec:primorial}] \textbf{Primorial deserts, made unconditional.} A
multiple of $\prod_{p\le P}p$ has no prime summand $\le P$
(Lemma~\ref{lem:primcore}). Combining this with Dirichlet
(Theorem~\ref{thm:dirichlet}) yields, for every prime $P$, an even $\N$ with an
\emph{explicit} partition and $\g(\N)=\mathrm{nextprime}(P)$; hence
$\sup_\N\g(\N)=\infty$ unconditionally, with $\g(\N)\gg\log \N$ by Linnik.
\item[\S\ref{sec:covering}] \textbf{Covering systems.} A greedy covering of the
odd primes below $Q$ forces ``no prime $<Q$ is a summand'' with modulus
$M\ll\prod_{p<Q}p$ (Proposition~\ref{prop:cover}); any covering can be realised
as an actual $\N$ with an explicit partition (Theorem~\ref{thm:coverdirichlet}),
so $\g(\N)\ge Q$ unconditionally. The size advantage over the primorial
\emph{widens} with $Q$ (Table~\ref{tab:results}).
\item[\S\ref{sec:gaps}] \textbf{Transfer from prime gaps.} A long prime gap
forces a large least summand; the converse fails. The resulting lower bound is
$\g(\N)\gg\frac{\log \N\,\log_2\N\,\log_4\N}{\log_3\N}$, which together
with~\eqref{eq:GLR} brackets the true maximal order.
\end{itemize}
We exhibit a verified $1020$-digit $\N$ with $\g(\N)=100\,747$ and certify a
$327\,455$-digit $\N$ with no prime summand below $10^8$
(\S\ref{sec:results}); a standalone verifier with hashes is in
Appendix~\ref{app:verify}.

\paragraph{Contribution.} We report two records, to the best of our knowledge.
The $\g(\N)=100\,747$ example is the largest \emph{proven} least Goldbach summand:
its prime complement $P$ is certified prime by an Atkin--Morain elliptic-curve
proof (Appendix~\ref{app:P}), so the value is unconditional and about ten times
the previous record of $9781$. The $Q=10^8$ row is the largest explicit
Goldbach-summand \emph{obstruction}: an integer carrying a short modular
certificate that rules out every prime offset below $10^8$, four orders of
magnitude beyond the reach of exhaustive search. We remain careful about status:
the modular certificates of absence are unconditional; the $100\,747$ value is
proven (deterministic for $100\,747$, ECPP for $P$); and the $10^8$ figure bounds
$\g(\N)$ from below only if $\N$ is a Goldbach number at all.

Throughout, $\log_k$ is the $k$-fold iterated logarithm (so $\log_2 \N=\log\log
\N$, $\log_3\N=\log\log\log \N$, and so on---each extra $\log$ grows
\emph{extremely} slowly); $\theta(x)=\sum_{p\le x}\log p$ is the sum of $\log p$
over all primes $p\le x$ (Chebyshev's function); $\gamma\approx0.5772$ is Euler's
constant; $f\gg g$ means $f\ge c\,g$ for some constant $c>0$ and all large
arguments (the reverse of $f=O(g)$); and ``BPSW'' denotes the Baillie--PSW
probable-prime test~\cite{PSW}.

\section{Factorial neighbourhoods and Wilson obstructions}\label{sec:factorial}

\begin{lemma}[Factorial composites]\label{lem:fact}
For $n\ge2$ and $2\le k\le n$, the number $n!+k$ is composite. For $n\ge4$ and
$2\le k\le n$, the number $n!-k$ is also composite.
\end{lemma}
\begin{proof}
In both cases $k\mid n!$, so $k\mid(n!\pm k)$. For the plus side $n!+k>k$, so
$n!+k$ is a proper multiple of $k$. For the minus side, $n!-k=k\bigl(n!/k-1\bigr)$
with $n!/k-1\ge2$ exactly when $n!\ge3k$, which holds for $n\ge4$; it fails at
$n=3,\,k=3$, where $3!-3=3$ is prime.
\end{proof}

\begin{lemma}[Wilson obstructions]\label{lem:wilson}
Let $n\ge2$. \emph{(A)} If $n+1$ is prime, then $(n+1)\mid n!+1$. \emph{(B)} If
$n+2$ is prime, then $(n+2)\mid n!-1$.
\end{lemma}
\noindent Recall \emph{Wilson's theorem}: a number $p>1$ is prime if and only if
$(p-1)!\equiv-1\pmod p$, i.e. the product $1\cdot2\cdots(p-1)$ is one less than a
multiple of $p$. We use only the forward direction.
\begin{proof}
For (A) put $p=n+1$; Wilson's theorem gives $n!=(p-1)!\equiv-1\pmod p$. For (B)
put $p=n+2$; then $(p-1)!=(n+1)\,n!$ and $n+1\equiv-1\pmod p$, so
$-n!\equiv-1$, i.e. $n!\equiv1\pmod p$.
\end{proof}

\begin{corollary}\label{cor:facprime}
If $n!+1>n+1$ is prime, then $n+1$ is composite. In particular the factorial
primes $n!+1$ (for $n=1,2,3,11,27,37,41,\dots$) sit at composite $n+1$.
\end{corollary}

\subsection{The worked example $\N=27!+28$}

Take $\N=27!+28$, even with $29$ digits. Since $27!\equiv0\pmod r$ for every prime
$r\le27$,
\begin{equation}\label{eq:shift}
\N-q\;=\;27!+(28-q)\;\equiv\;28-q\pmod r\qquad(r\le27\ \text{prime}),
\end{equation}
so $\N-q$ is killed by a small prime $r$ exactly when $r\mid(28-q)$: a shifted
sieve that murders the class $q\equiv28\pmod r$.

\begin{proposition}\label{prop:71}
For $\N=27!+28$ every prime $q$ with $2\le q\le67$ has $\N-q$ composite, while
$\N-71$ is prime; hence $\g(\N)=71$.
\end{proposition}
\begin{proof}[Proof by explicit witnesses]
Table~\ref{tab:witness} gives, for each prime $q\le71$, a factor of $\N-q$. For
$q\le23$ the witness divides $28-q$ via~\eqref{eq:shift}. Two cases below the
winner are notable: $q=29$ is killed by $29$ because $29=27+2$ is prime, so
Lemma~\ref{lem:wilson}(B) gives $29\mid27!-1=\N-29$; and $q=59$ survives the
sieve (as $59-28=31$ is $27$-rough) yet $41\mid27!-31=\N-59$. Finally
$71-28=43$ is $27$-rough and $27!-43$ is prime.
\end{proof}

\begin{table}[h]
\centering\small
\caption{Why every small prime fails as a summand of $\N=27!+28$.}\label{tab:witness}
\begin{tabular}{rlc@{\qquad}rlc}
\toprule
$q$ & $\N-q$ & factor & $q$ & $\N-q$ & factor\\
\midrule
$2$&$27!+26$&$2$ & $31$&$27!-3$&$3$\\
$3$&$27!+25$&$5$ & $37$&$27!-9$&$3$\\
$5$&$27!+23$&$23$ & $41$&$27!-13$&$13$\\
$7$&$27!+21$&$3$ & $43$&$27!-15$&$3$\\
$11$&$27!+17$&$17$ & $47$&$27!-19$&$19$\\
$13$&$27!+15$&$3$ & $53$&$27!-25$&$5$\\
$17$&$27!+11$&$11$ & $59$&$27!-31$&$\mathbf{41}$ (generic)\\
$19$&$27!+9$&$3$ & $61$&$27!-33$&$3$\\
$23$&$27!+5$&$5$ & $67$&$27!-39$&$3$\\
$29$&$27!-1$&$\mathbf{29}$ (Wilson) & $71$&$27!-43$&\textbf{prime}\\
\bottomrule
\end{tabular}
\end{table}

By~\eqref{eq:shift} a prime $q$ can even be \emph{tested} only if $q-28$ is
$27$-rough, so the viable candidates $29,59,71,89,\dots$ are sparse, and $71$ is
merely the first with a prime complement. Sharper still: $27!+1$ \emph{is} prime
(Corollary~\ref{cor:facprime}), the natural bridge $q=27$, but $27$ is composite,
so one falls through to $71$. For indices that are simultaneously prime and
factorial-prime indices ($n=3,11,37,41$) the bridge works and
$\g(n!+(n+1))=n$. On its own the factorial construction reaches only
$\g\approx2n\approx2\log\N/\log\log\N$; the next sections do far better.

\section{Primorial deserts, made unconditional}\label{sec:primorial}

The factorial wastes each small prime on a residue class. Taking $\N$ to be a
multiple of a primorial wastes nothing.

\begin{lemma}[Primorial core]\label{lem:primcore}
Let $P$ be prime, let $m_P=\prod_{p\le P}p$ be the \emph{primorial}---the product
of every prime up to $P$ (for instance $m_{11}=2\cdot3\cdot5\cdot7\cdot11=2310$)
---and let $\N$ be a multiple of $m_P$ with $\N>2P$. Then $\N-q$ is composite for
every prime $q\le P$; i.e. no prime $\le P$ is a Goldbach summand of $\N$.
\end{lemma}
\begin{proof}
For a prime $q\le P$ we have $q\mid m_P\mid\N$ (as $q$ is one of the factors of
$m_P$), hence $q\mid\N-q$, and $\N-q>P\ge q$, so $\N-q$ is a proper multiple of
$q$ and therefore composite.
\end{proof}

Lemma~\ref{lem:primcore} alone gives only claim~(a) of the introduction. To reach
unboundedness with an \emph{actual} Goldbach number we add Dirichlet's theorem.

\begin{theorem}[Unconditional unboundedness]\label{thm:dirichlet}
For every prime $P$ put $s=\mathrm{nextprime}(P)$. There are infinitely many
primes $t\equiv-s\pmod{m_P}$, and for any such $t$ the even integer $\N=s+t$
satisfies $m_P\mid\N$, has the explicit Goldbach partition $\N=s+t$, and has
\[
\g(\N)\;=\;\mathrm{nextprime}(P)\;>\;P.
\]
Hence $\sup_\N\g(\N)=\infty$ unconditionally. Moreover the least such $t$
satisfies $\log\N\ll_{\,}P$ by Linnik's theorem~\cite{Linnik}, so
$\g(\N)\gg\log \N$.
\end{theorem}
\begin{proof}
As $\gcd(s,m_P)=1$ (because $s>P$), Dirichlet~\cite{Dirichlet} gives infinitely
many primes $t\equiv-s\pmod{m_P}$. Then $\N=s+t\equiv0\pmod{m_P}$ and $\N$ is
even. By Lemma~\ref{lem:primcore} no prime $\le P$ is a summand; the only primes
below $s=\mathrm{nextprime}(P)$ are those $\le P$, so the least summand is at
least $s$. But $\N-s=t$ is prime, so $s$ is a summand; hence $\g(\N)=s$. Linnik's
theorem bounds the least such prime by $t\ll m_P^{L}$ for an absolute constant
$L$, and $\log m_P=\theta(P)\sim P$, so $\log\N\ll P\sim s=\g(\N)$.
\end{proof}

\begin{remark}[Concrete instances]\label{rem:concrete}
For $P=13$ one has $m_{13}=30030$, $s=17$, and the least prime $t\equiv-17$
$\pmod{30030}$ is $t=30013$, giving $\N=30030=2\cdot3\cdot5\cdot7\cdot11\cdot13$
with partition $30030=17+30013$ and $\g(\N)=17$---all primality deterministic. A
larger fully proven example is the primorial $m_{43}=13082761331670030$ itself:
its minimal partition is $m_{43}=89+13082761331669941$ (both proven prime, being
$<2^{64}$), so $\g(m_{43})=89$ unconditionally.
\end{remark}

\begin{remark}[The roughness bonus]\label{rem:rough}
For the bare choice $\N=m_P$, a prime $q>P$ makes $\N-q$ coprime to all primes
$\le P$---a property called \emph{$P$-rough}---which improves its chances of being
prime, since the small primes that catch most composites are already excluded. By
\emph{Mertens' theorem} (which says $\prod_{p\le P}(1-1/p)\sim e^{-\gamma}/\log P$,
the product of $1-1/p$ over primes $p\le P$ shrinking like $1/\log P$), this
sieving effect multiplies the heuristic probability by
$\prod_{p\le P}(1-1/p)^{-1}\sim e^\gamma\log P$. The expected least surviving
summand then lands near $P(1+e^{-\gamma})\approx1.56\,P$. Either way the primorial
pays a modulus $m_P\approx e^{P}$ (its size is $\log m_P=\theta(P)\sim P$, so $m_P$
itself is exponentially large in $P$); the next section removes this waste.
\end{remark}

\section{Covering systems: large deserts cheaply}\label{sec:covering}

\subsection{The construction}

A primorial spends one prime per summand; a covering system reuses each prime on
a whole residue class. The one structural constraint is parity: a sum of two odd
primes is even, so $\N$ is kept even and the prime $2$ kills only $q=2$; every
\emph{odd} prime below the target must be covered by odd sieving primes.

\begin{proposition}[Covering certificate]\label{prop:cover}
Let $Q\ge3$ and let $\{(r,a_r):r\in\mathcal R\}$ be a system of odd primes
$\mathcal R$ such that every odd prime $q<Q$ has $q\equiv a_r\pmod r$ for some
$r\in\mathcal R$. If $\N$ is even, $\N\equiv a_r\pmod r$ for all $r\in\mathcal R$,
and
\[
\N\;>\;Q+\max_{r\in\mathcal R}r,
\]
then $\N-q$ is composite for every prime $q<Q$; i.e. no prime $<Q$ is a Goldbach
summand of $\N$.
\end{proposition}
\begin{proof}
For $q=2$, $\N-2$ is even and $>2$. For an odd prime $q<Q$, pick
$r\in\mathcal R$ with $q\equiv a_r\equiv\N\pmod r$; then $r\mid\N-q$, and
$\N-q>(Q+\max_r r)-Q\ge r$, so $r\mid\N-q$ properly and $\N-q$ is composite.
\end{proof}

\begin{remark}
The size hypothesis only has to make every divisor $r\mid\N-q$ proper, for which
$\N>Q+\max_r r$ suffices; the much stronger $\N>Q\prod_r r$ is unnecessary. A
common solution $\N$ exists modulo $M=2\prod_{r}r$ by the Chinese Remainder
Theorem, and any large enough even representative qualifies.
\end{remark}

We build the covering greedily and deterministically.

\medskip
\noindent\fbox{\parbox{0.97\textwidth}{%
\textbf{Algorithm (greedy covering of the odd primes below $Q$).}
\begin{enumerate}\setlength\itemsep{1pt}
\item Let $T=\{\text{odd primes }<Q\}$ and require $\N\equiv0\pmod2$.
\item For $r=3,5,7,11,\dots$ in increasing order, while $T\ne\varnothing$: let
$a_r$ be a residue mod $r$ containing the most elements of $T$ (smallest such
residue on ties), require $\N\equiv a_r\pmod r$, and delete from $T$ every
$q\equiv a_r\pmod r$.
\item Recover $\N$ from the congruences by the Chinese Remainder Theorem.
\end{enumerate}}}
\medskip

The resulting modulus is $M=2\prod_{r\in\mathcal R}r$---twice the product of the
sieving primes actually used---which is far smaller than the primorial
$\prod_{p<Q}p$ of all primes below $Q$, precisely because each sieving prime $r$
is reused to kill an entire residue class of targets rather than a single one.

\subsection{Realising a covering unconditionally}

Proposition~\ref{prop:cover} still only gives claim~(a). The same Dirichlet idea
as in \S\ref{sec:primorial}, applied twice, upgrades any covering to an actual
Goldbach number.

\begin{theorem}[Unconditional realisation]\label{thm:coverdirichlet}
Let $Q\ge3$ and let $\{(r,a_r):r\in\mathcal R\}$ be a covering of the odd primes
below $Q$ by odd primes $\mathcal R$. Then there are infinitely many even
integers $\N$ admitting an explicit Goldbach partition $\N=s+t$ ($s,t$ prime)
with $\N\equiv a_r\pmod r$ for all $r\in\mathcal R$. For every such $\N$, no prime
below $Q$ is a summand, so $\g(\N)$ exists and $\g(\N)\ge Q$.
\end{theorem}
\begin{proof}
For each $r\in\mathcal R$ choose $b_r\in\{1,\dots,r-1\}\setminus\{a_r\}$, possible
since $r\ge3$. By CRT and Dirichlet there is a prime $s$ with $s\equiv b_r\pmod r$
for all $r$; as $s\not\equiv a_r\pmod r$ for any $r$ and every odd prime $<Q$ is
covered, necessarily $s\ge Q$. Next require $t\equiv a_r-s\pmod r$ for all $r$ and
$t\equiv1\pmod2$. Each $a_r-s\not\equiv0\pmod r$, so the combined modulus
$2\prod_r r$ and residue $c$ satisfy $\gcd(c,2\prod_r r)=1$; Dirichlet gives
infinitely many primes $t\equiv c$. For any such $t$, set $\N=s+t$: then
$\N\equiv a_r\pmod r$ for all $r$, $\N$ is even, and $\N=s+t$ is an explicit
partition. For a prime $q<Q$: if $q=2$ then $\N-2$ is even and composite; if $q$
is odd then $q\equiv a_r\equiv\N\pmod r$ for some $r$, and for $\N$ large enough
$r\mid\N-q$ properly. Thus no prime $<Q$ is a summand and $\g(\N)\ge Q$.
\end{proof}

\subsection{Transfer from large prime gaps}\label{sec:gaps}

\begin{proposition}\label{prop:transfer}
If $[\N-Q,\N-2]$ contains no prime, then $\g(\N)>Q$.
\end{proposition}
\begin{proof}
Every prime $q\in[2,Q]$ has $\N-q\in[\N-Q,\N-2]$, hence composite.
\end{proof}

So a prime gap of length $\ge Q$ just below $\N$ forces $\g(\N)>Q$: \emph{prime
gaps transfer to least Goldbach summands}. The transfer is, however, strictly
one-directional. Forcing $\g(\N)>Q$ requires $\N-q$ composite only for
\emph{prime} offsets $q<Q$, not for every integer in an interval; the covering of
Proposition~\ref{prop:cover} covers just the primes below $Q$, a far sparser
target. Thus ``large least Goldbach summand'' is a relaxation of ``large prime
gap,'' and one should not expect the two problems to coincide.

Writing $G(X)$ for the largest gap between consecutive primes below $X$, the work
of Ford--Green--Konyagin--Tao and of Maynard~\cite{FGKT,Maynard}, sharpened
in~\cite{FGKMT}, gives
\[
G(X)\;\gg\;\frac{\log X\,\log_2 X\,\log_4 X}{\log_3 X},
\]
improving Rankin~\cite{Rankin}. In words: there are stretches of consecutive
composite integers below $X$ of length at least a constant times
$\log X\cdot\log_2 X\cdot\log_4 X/\log_3 X$. The dominant factor is $\log X$ (the
average prime gap near $X$); the three iterated-log factors are a slowly growing
extra boost, and since each $\log_k$ grows glacially this is only a little more
than $\log X$. Feeding the covering that underlies these bounds into
Theorem~\ref{thm:coverdirichlet} (which costs only a constant factor in
$\log\N$, by Linnik) makes the transfer unconditional:

\begin{theorem}\label{thm:lower}
There are infinitely many even integers $\N$, each with an explicit Goldbach
representation, such that
\[
\g(\N)\;\gg\;\frac{\log \N\,\log_2 \N\,\log_4 \N}{\log_3 \N}.
\]
\end{theorem}

Together with the Granville--van de Lune--te Riele
conjecture~\eqref{eq:GLR}, $\g(\N)=O((\log \N)^2\log\log \N)$, this brackets the
maximal order of $\g$ between a logarithmic-type lower bound and a
$(\log \N)^2$-type conjectural ceiling.

\section{Computational results}\label{sec:results}

\paragraph{Method.} The greedy covering was run with a vectorised sieve and $\N$
assembled by the Chinese Remainder Theorem; the construction is deterministic, so
fixing the prime ordering and the tie rule reproduces every number below.
Candidate complements were screened with Baillie--PSW, which is a
\emph{probable}-prime test rather than a primality certificate~\cite{PSW}; the
headline complement was then \emph{proven} prime by an Atkin--Morain
elliptic-curve (ECPP) certificate~\cite{AtkinMorain} (the criterion is stated in
Appendix~\ref{app:P}). Complements below $2^{64}$ (and the small examples of
Remark~\ref{rem:concrete}) are settled by a deterministic test.

\paragraph{Certified, conditional, and proven.} Three statuses appear in
Table~\ref{tab:results}. \emph{Certified} (unconditional, no primality test on
$\N$): for each $Q$ we checked exhaustively that the covering kills every odd
prime below $Q$, so no prime $<Q$ is a summand. \emph{Conditional}: for $Q=10^4$
the least surviving complement is only a BPSW probable prime, so $\g(\N)=10\,159$
rests on its (unproven) primality. \emph{Proven}: for $Q=10^5$ the complement
$P=\N-100\,747$ carries an ECPP certificate (Appendix~\ref{app:P}), so
$\g(\N)=100\,747$ is unconditional, as is the small example $\g(m_{43})=89$ of
Remark~\ref{rem:concrete}.

\begin{table}[h]
\centering
\caption{Forcing the least Goldbach summand. ``Congruences'' is
$|\mathcal R|+1$, including the parity congruence $2\,{:}\,0$; ``Primorial
digits'' is $\log_{10}\prod_{p\le \text{gap}}p$, the size the primorial of
Lemma~\ref{lem:primcore} would need for the same gap; ``Killed/$r$'' is the
average number of odd primes removed per odd sieving prime; ``Shrink'' is the
ratio of the two size columns. $^\dagger$Conditional on the BPSW
probable-primality of the complement $\N-\g$; $^\ddagger$proven, the complement
carrying an ECPP certificate. For $Q\ge10^6$ only the
certified statement (no prime $<Q$ is a summand) is claimed.}\label{tab:results}
\begin{tabular}{rrrrrrrr}
\toprule
$Q$ & $\g(\N)$ & Congr. & Largest & Digits & Primorial & Killed/ & Shrink\\
 & & $|\mathcal R|{+}1$ & $r$ & $\N$ & digits & $r$ & \\
\midrule
$10^{4}$ & $10{,}159^\dagger$ & $86$ & $443$ & $182$ & $4{,}370$ & $14.4$ & $24.0\times$\\
$10^{5}$ & $100{,}747^\ddagger$ & $356$ & $2{,}393$ & $1{,}020$ & $43{,}603$ & $27.0$ & $42.7\times$\\
$10^{6}$ & no prime $<10^{6}$ & $1{,}762$ & $15{,}083$ & $6{,}480$ & $433{,}636$ & $44.6$ & $66.9\times$\\
$10^{7}$ & no prime $<10^{7}$ & $9{,}853$ & $102{,}983$ & $44{,}599$ & $4{,}340{,}851$ & $67.5$ & $97.3\times$\\
$10^{8}$ & no prime $<10^{8}$ & $60{,}627$ & $755{,}137$ & $327{,}455$ & $43{,}424{,}120$ & $95.0$ & $132.6\times$\\
\bottomrule
\end{tabular}
\end{table}

\paragraph{The exhibited number.} The row $Q=10^{5}$ gives a $1020$-digit even
integer $\N$ (Appendix~\ref{app:N}) for which $\N-q$ is composite for all $9592$
primes $q<10^{5}$. Its least surviving prime is $100\,747$, with
$P=\N-100\,747$ a $1020$-digit prime, proved prime by an Atkin--Morain ECPP
certificate (Appendix~\ref{app:P}). Since $100\,747$ is (deterministically) prime
and the partition $\N=100\,747+P$ is exhibited, $\g(\N)=100\,747$
\emph{unconditionally}. The defining
congruence system is in Appendix~\ref{app:res}; the SHA-256 hashes are
\[
\texttt{N}:\ \texttt{853e0013\dots e729bfc},\quad
\texttt{P}:\ \texttt{32894e4b\dots 6e367838},\quad
\text{ECPP cert}:\ \texttt{9ebadbc4\dots 41edafe},
\]
and a standalone verifier is in Appendix~\ref{app:verify}.

\paragraph{The certified number.} The row $Q=10^{8}$ defines an even integer with
$\log_{10}M=327{,}455$ for which $\N-q$ is composite for all $5{,}761{,}455$
primes $q<10^{8}$, checked against the covering. This certifies, unconditionally,
that no prime below $10^{8}$ is a Goldbach summand; whether $\N$ has any summand
at all is conditional on Goldbach's conjecture. Unconditionally,
Theorem~\ref{thm:coverdirichlet} guarantees a (larger) even integer with an
explicit partition and $\g\ge10^{8}$.

\paragraph{Reading the table.} ``Killed/$r$'' rises from $14$ to $95$: each
sieving prime is reused ever more heavily, the advantage a covering has over the
one-for-one primorial. ``Shrink'' rises from $24\times$ to $133\times$: that
advantage \emph{widens} with $Q$. The gap of $10^{8}$ costing the covering
$327{,}455$ digits would cost the primorial more than $43$ million.

\section{Discussion and open questions}

We have driven the least Goldbach summand from an incidental $71$ near $27!$ to a
verified $100\,747$ and a certified $10^{8}$, with the unboundedness statement
made unconditional through Dirichlet. The picture that emerges is a sandwich:
\[
\underbrace{\frac{\log \N\,\log_2 \N\,\log_4 \N}{\log_3 \N}}_{\text{rigorous lower bound (Thm.~\ref{thm:lower})}}
\;\ll\;\max_{\N'\le \N}\g(\N')\;\lesssim\;\underbrace{(\log \N)^2\log\log \N}_{\text{conjectural (GLvdtR)}}.
\]
Our explicit greedy constructions sit well below even the lower-bound rate,
because a greedy covering is far from the optimal Erd\H{o}s--Rankin one; they are
nonetheless enough to demonstrate unboundedness with concrete, certified numbers.

It is worth separating two distinct goals. One is to make $\g(\N)$ \emph{large},
disregarding the size of $\N$; here the covering method excels and scales to
obstructions ($10^8$) far beyond the reach of any exhaustive search. The other is
the smallest-desert problem $S(q)$---to make $\g(\N)=q$ with $\N$ as \emph{small}
as possible; here the method is highly suboptimal. The exhaustive record
$\g=9781$ sits at a $19$-digit integer, whereas our comparable $\g=10{,}159$
needs a $182$-digit CRT integer. The value of the greedy covering is therefore
not efficiency but transparency: it produces a short, independently checkable
modular certificate at any prescribed scale, where exhaustive search cannot
reach and where the optimal $S(q)$ remains out of analytic grasp.

\begin{enumerate}
\item \textbf{The smallest desert $S(q)$.} What is the least even $\N$ with
$\g(\N)=q$~\cite{GLvdtR}? This is the natural extremal companion to the largest
prime gap, and the greedy moduli here are only upper bounds on it.
\item \textbf{Removing the conditional.} The $10^5$ desert is now fully
unconditional, its complement ECPP-certified. The $10^8$ certificate is a
different matter: ruling out small summands is unconditional, but exhibiting
\emph{any} partition of that $327{,}455$-digit integer needs a prime $q\ge10^{8}$
with $\N-q$ prime, then a primality proof of that $\sim\!10^{8}$-times-larger
complement---out of reach at present.
\item \textbf{Sparser than gaps.} Does the restriction to prime offsets let
$\g(\N)$ exceed the maximal-prime-gap rate, closing some of the distance to the
conjectural ceiling~\eqref{eq:GLR}?
\end{enumerate}

The moral is a trade-off, not a paradox: $\g(\N)$ can be made as large as one
likes, but only by paying in the size of $\N$, at an exchange rate governed by
the growth of gaps between primes.

\begin{thebibliography}{99}
\bibitem{OeS} T.~Oliveira e Silva, S.~Herzog, S.~Pardi, \emph{Empirical
verification of the even Goldbach conjecture and computation of prime gaps up to
$4\cdot10^{18}$}, Math. Comp. \textbf{83} (2014), 2033--2060.
\bibitem{Richstein} J.~Richstein, \emph{Verifying the Goldbach conjecture up to
$4\cdot10^{14}$}, Math. Comp. \textbf{70} (2001), 1745--1749.
\bibitem{HL} G.~H.~Hardy, J.~E.~Littlewood, \emph{Some problems of `Partitio
Numerorum' III}, Acta Math. \textbf{44} (1923), 1--70.
\bibitem{GLvdtR} A.~Granville, J.~van de Lune, H.~J.~J.~te Riele, \emph{Checking
the Goldbach conjecture on a vector computer}, in: Number Theory and Applications
(R.~A.~Mollin, ed.), NATO ASI Ser.~C \textbf{265}, Kluwer, Dordrecht, 1989,
423--433.
\bibitem{OEIS} N.~J.~A.~Sloane (ed.), \emph{The On-Line Encyclopedia of Integer
Sequences}, sequences A020481 (least Goldbach summand) and A025018/A025019
(records), \url{https://oeis.org}.
\bibitem{Dirichlet} P.~G.~L.~Dirichlet, \emph{Beweis des Satzes, dass jede
unbegrenzte arithmetische Progression\dots\ unendlich viele Primzahlen
enth\"alt}, Abh. K. Preuss. Akad. Wiss. (1837), 45--81.
\bibitem{Linnik} U.~V.~Linnik, \emph{On the least prime in an arithmetic
progression I, II}, Mat. Sbornik \textbf{15(57)} (1944), 139--178, 347--368.
\bibitem{Rankin} R.~A.~Rankin, \emph{The difference between consecutive prime
numbers}, J. London Math. Soc. \textbf{13} (1938), 242--247.
\bibitem{FGKT} K.~Ford, B.~Green, S.~Konyagin, T.~Tao, \emph{Large gaps between
consecutive prime numbers}, Ann. of Math. \textbf{183} (2016), 935--974.
\bibitem{Maynard} J.~Maynard, \emph{Large gaps between primes}, Ann. of Math.
\textbf{183} (2016), 915--933.
\bibitem{FGKMT} K.~Ford, B.~Green, S.~Konyagin, J.~Maynard, T.~Tao, \emph{Long
gaps between primes}, J. Amer. Math. Soc. \textbf{31} (2018), 65--105.
\bibitem{PSW} C.~Pomerance, J.~L.~Selfridge, S.~S.~Wagstaff~Jr., \emph{The
pseudoprimes to $25\cdot10^9$}, Math. Comp. \textbf{35} (1980), 1003--1026.
\bibitem{AtkinMorain} A.~O.~L.~Atkin, F.~Morain, \emph{Elliptic curves and
primality proving}, Math. Comp. \textbf{61} (1993), 29--68.
\bibitem{Cramer} H.~Cram\'er, \emph{On the order of magnitude of the difference
between consecutive prime numbers}, Acta Arith. \textbf{2} (1936), 23--46.
\end{thebibliography}

\appendix

\section{The verified integer $\N$ ($1020$ digits)}\label{app:N}
Even, with $\g(\N)=100\,747$ (unconditional; the complement $P$ is ECPP-certified,
Appendix~\ref{app:P}). It is the Chinese Remainder Theorem solution of the system
of Appendix~\ref{app:res}, reduced to its least representative exceeding $10^{6}$.
\begin{verbatim}
1069544539282386715486223835075605584373382706541914908831505623262975
2083522612338734637321305998987592550408498828612570022394208940143460
0064975415094707742731050071929426768009802994798394217849911437581801
4039594735967522377970546037935637237928281016918193304430368169807038
4474941987767969933359820683681557599943056143302527468556049195060358
4280425970311500284306417878555994564603068444664362599417258644737229
0484754637092494892808114334223748793668446381738325290148115072022495
7522077467401451931244937454496416229136154184084575547920258279665074
5524835562942450136969422849093934585977269428371274331010556593344957
8955544532422815286226956359279490135931859415743367621432745233231595
6982087156900220665027823763411480472052164539053448304868021268792102
4743632536289493802245261747302618432579314092247670538805639533124414
6693007026522943169287884632030399873821024301289446160936362452780885
7247690680844939813744595253510987977660444705920124097262087259987624
2505737637074064752090480111959678911758
\end{verbatim}


\section{The complement $P=\N-100\,747$ ($1020$ digits)}\label{app:P}
A prime, proved by an Atkin--Morain elliptic-curve certificate~\cite{AtkinMorain}
(PARI \texttt{primecert}; SHA-256 \texttt{9ebadbc4\dots41edafe}). The certificate
is a chain of steps, each invoking the following criterion.

\smallskip
\noindent\textbf{Elliptic-curve primality criterion (Goldwasser--Kilian).} Let
$N>1$ be coprime to $6$, let $E:\,y^2=x^3+ax+b$ be an elliptic curve over
$\mathbb Z/N\mathbb Z$, and let $m,q$ be integers with $q$ prime, $q\mid m$, and
$q>(N^{1/4}+1)^2$. If there is a point $\mathbf P$ on $E$ with $m\,\mathbf P=O$
(the point at infinity) while $(m/q)\,\mathbf P\neq O$, then $N$ is prime.

\smallskip
\noindent Here $m$ is the number of points on $E$ and $q$ is a large prime factor
of it; the bound $q>(N^{1/4}+1)^2$ is what makes the implication go through. Each
step thus reduces the primality of $N$ to that of the strictly smaller prime $q$,
which becomes the $N$ of the next step (the ``downrun''); after $120$ steps the
chain falls from $1020$ to $16$ digits, where the final $q$ is small enough to
prove directly. With $100\,747$ this gives the partition $\N=100\,747+P$, so
$\g(\N)=100\,747$ unconditionally.
\begin{verbatim}
1069544539282386715486223835075605584373382706541914908831505623262975
2083522612338734637321305998987592550408498828612570022394208940143460
0064975415094707742731050071929426768009802994798394217849911437581801
4039594735967522377970546037935637237928281016918193304430368169807038
4474941987767969933359820683681557599943056143302527468556049195060358
4280425970311500284306417878555994564603068444664362599417258644737229
0484754637092494892808114334223748793668446381738325290148115072022495
7522077467401451931244937454496416229136154184084575547920258279665074
5524835562942450136969422849093934585977269428371274331010556593344957
8955544532422815286226956359279490135931859415743367621432745233231595
6982087156900220665027823763411480472052164539053448304868021268792102
4743632536289493802245261747302618432579314092247670538805639533124414
6693007026522943169287884632030399873821024301289446160936362452780885
7247690680844939813744595253510987977660444705920124097262087259987624
2505737637074064752090480111959678811011
\end{verbatim}


\section{The covering congruence system}\label{app:res}
Each entry $r\,{:}\,a$ means $\N\equiv a\pmod r$. The first, $2\,{:}\,0$, fixes
parity; the remaining $355$ odd primes (so $|\mathcal R|=355$, and $356$
congruences in all) cover every odd prime below $10^{5}$. Applying the Chinese
Remainder Theorem and taking the least representative above $10^{6}$ reproduces
Appendix~\ref{app:N}.
\begin{center}\scriptsize
\setlength{\tabcolsep}{4pt}
\begin{tabular}{rrrrrrrr}
\toprule
$2\,{:}\,0$ & $3\,{:}\,2$ & $5\,{:}\,3$ & $7\,{:}\,5$ & $11\,{:}\,8$ & $13\,{:}\,6$ & $17\,{:}\,4$ & $19\,{:}\,11$ \\
$23\,{:}\,10$ & $29\,{:}\,7$ & $31\,{:}\,17$ & $37\,{:}\,20$ & $41\,{:}\,19$ & $43\,{:}\,12$ & $47\,{:}\,37$ & $53\,{:}\,35$ \\
$59\,{:}\,37$ & $61\,{:}\,22$ & $67\,{:}\,11$ & $71\,{:}\,6$ & $73\,{:}\,60$ & $79\,{:}\,43$ & $83\,{:}\,6$ & $89\,{:}\,59$ \\
$97\,{:}\,7$ & $101\,{:}\,75$ & $103\,{:}\,83$ & $107\,{:}\,33$ & $109\,{:}\,87$ & $113\,{:}\,24$ & $127\,{:}\,53$ & $131\,{:}\,57$ \\
$137\,{:}\,2$ & $139\,{:}\,31$ & $149\,{:}\,51$ & $151\,{:}\,61$ & $157\,{:}\,53$ & $163\,{:}\,31$ & $167\,{:}\,124$ & $173\,{:}\,33$ \\
$179\,{:}\,40$ & $181\,{:}\,162$ & $191\,{:}\,62$ & $193\,{:}\,109$ & $197\,{:}\,51$ & $199\,{:}\,33$ & $211\,{:}\,30$ & $223\,{:}\,209$ \\
$227\,{:}\,18$ & $229\,{:}\,32$ & $233\,{:}\,123$ & $239\,{:}\,152$ & $241\,{:}\,3$ & $251\,{:}\,42$ & $257\,{:}\,185$ & $263\,{:}\,139$ \\
$269\,{:}\,18$ & $271\,{:}\,21$ & $277\,{:}\,211$ & $281\,{:}\,242$ & $283\,{:}\,235$ & $293\,{:}\,101$ & $307\,{:}\,30$ & $311\,{:}\,54$ \\
$313\,{:}\,150$ & $317\,{:}\,15$ & $331\,{:}\,199$ & $337\,{:}\,145$ & $347\,{:}\,250$ & $349\,{:}\,125$ & $353\,{:}\,29$ & $359\,{:}\,47$ \\
$367\,{:}\,101$ & $373\,{:}\,288$ & $379\,{:}\,82$ & $383\,{:}\,338$ & $389\,{:}\,319$ & $397\,{:}\,83$ & $401\,{:}\,106$ & $409\,{:}\,318$ \\
$419\,{:}\,64$ & $421\,{:}\,143$ & $431\,{:}\,15$ & $433\,{:}\,124$ & $439\,{:}\,3$ & $443\,{:}\,27$ & $449\,{:}\,231$ & $457\,{:}\,249$ \\
$461\,{:}\,41$ & $463\,{:}\,142$ & $467\,{:}\,82$ & $479\,{:}\,3$ & $487\,{:}\,104$ & $491\,{:}\,33$ & $499\,{:}\,243$ & $503\,{:}\,51$ \\
$509\,{:}\,165$ & $521\,{:}\,214$ & $523\,{:}\,496$ & $541\,{:}\,39$ & $547\,{:}\,201$ & $557\,{:}\,518$ & $563\,{:}\,548$ & $569\,{:}\,446$ \\
$571\,{:}\,127$ & $577\,{:}\,499$ & $587\,{:}\,45$ & $593\,{:}\,20$ & $599\,{:}\,463$ & $601\,{:}\,78$ & $607\,{:}\,402$ & $613\,{:}\,305$ \\
$617\,{:}\,182$ & $619\,{:}\,16$ & $631\,{:}\,42$ & $641\,{:}\,43$ & $643\,{:}\,302$ & $647\,{:}\,7$ & $653\,{:}\,283$ & $659\,{:}\,479$ \\
$661\,{:}\,60$ & $673\,{:}\,198$ & $677\,{:}\,67$ & $683\,{:}\,500$ & $691\,{:}\,155$ & $701\,{:}\,46$ & $709\,{:}\,42$ & $719\,{:}\,152$ \\
$727\,{:}\,334$ & $733\,{:}\,284$ & $739\,{:}\,282$ & $743\,{:}\,257$ & $751\,{:}\,70$ & $757\,{:}\,364$ & $761\,{:}\,541$ & $769\,{:}\,542$ \\
$773\,{:}\,36$ & $787\,{:}\,28$ & $797\,{:}\,182$ & $809\,{:}\,7$ & $811\,{:}\,704$ & $821\,{:}\,506$ & $823\,{:}\,802$ & $827\,{:}\,86$ \\
$829\,{:}\,580$ & $839\,{:}\,4$ & $853\,{:}\,173$ & $857\,{:}\,91$ & $859\,{:}\,398$ & $863\,{:}\,287$ & $877\,{:}\,216$ & $881\,{:}\,14$ \\
$883\,{:}\,112$ & $887\,{:}\,123$ & $907\,{:}\,283$ & $911\,{:}\,8$ & $919\,{:}\,235$ & $929\,{:}\,249$ & $937\,{:}\,585$ & $941\,{:}\,96$ \\
$947\,{:}\,10$ & $953\,{:}\,347$ & $967\,{:}\,950$ & $971\,{:}\,68$ & $977\,{:}\,551$ & $983\,{:}\,179$ & $991\,{:}\,638$ & $997\,{:}\,71$ \\
$1009\,{:}\,482$ & $1013\,{:}\,2$ & $1019\,{:}\,113$ & $1021\,{:}\,880$ & $1031\,{:}\,592$ & $1033\,{:}\,302$ & $1039\,{:}\,352$ & $1049\,{:}\,803$ \\
$1051\,{:}\,495$ & $1061\,{:}\,865$ & $1063\,{:}\,359$ & $1069\,{:}\,216$ & $1087\,{:}\,851$ & $1091\,{:}\,199$ & $1093\,{:}\,602$ & $1097\,{:}\,748$ \\
$1103\,{:}\,926$ & $1109\,{:}\,305$ & $1117\,{:}\,692$ & $1123\,{:}\,457$ & $1129\,{:}\,13$ & $1151\,{:}\,489$ & $1153\,{:}\,694$ & $1163\,{:}\,34$ \\
$1171\,{:}\,347$ & $1181\,{:}\,61$ & $1187\,{:}\,362$ & $1193\,{:}\,79$ & $1201\,{:}\,113$ & $1213\,{:}\,660$ & $1217\,{:}\,1121$ & $1223\,{:}\,85$ \\
$1229\,{:}\,331$ & $1231\,{:}\,118$ & $1237\,{:}\,779$ & $1249\,{:}\,21$ & $1259\,{:}\,402$ & $1277\,{:}\,944$ & $1279\,{:}\,47$ & $1283\,{:}\,1138$ \\
$1289\,{:}\,761$ & $1291\,{:}\,292$ & $1297\,{:}\,34$ & $1301\,{:}\,625$ & $1303\,{:}\,1076$ & $1307\,{:}\,843$ & $1319\,{:}\,116$ & $1321\,{:}\,370$ \\
$1327\,{:}\,209$ & $1361\,{:}\,18$ & $1367\,{:}\,63$ & $1373\,{:}\,333$ & $1381\,{:}\,555$ & $1399\,{:}\,77$ & $1409\,{:}\,892$ & $1423\,{:}\,112$ \\
$1427\,{:}\,82$ & $1429\,{:}\,19$ & $1433\,{:}\,996$ & $1439\,{:}\,27$ & $1447\,{:}\,862$ & $1451\,{:}\,242$ & $1453\,{:}\,439$ & $1459\,{:}\,217$ \\
$1471\,{:}\,148$ & $1481\,{:}\,123$ & $1483\,{:}\,243$ & $1487\,{:}\,709$ & $1489\,{:}\,175$ & $1493\,{:}\,4$ & $1499\,{:}\,848$ & $1511\,{:}\,308$ \\
$1523\,{:}\,267$ & $1531\,{:}\,479$ & $1543\,{:}\,20$ & $1549\,{:}\,43$ & $1553\,{:}\,445$ & $1559\,{:}\,20$ & $1567\,{:}\,603$ & $1571\,{:}\,1$ \\
$1579\,{:}\,89$ & $1583\,{:}\,1466$ & $1597\,{:}\,1133$ & $1601\,{:}\,1272$ & $1607\,{:}\,7$ & $1609\,{:}\,163$ & $1613\,{:}\,1117$ & $1619\,{:}\,1100$ \\
$1621\,{:}\,802$ & $1627\,{:}\,547$ & $1637\,{:}\,250$ & $1657\,{:}\,702$ & $1663\,{:}\,142$ & $1667\,{:}\,1045$ & $1669\,{:}\,55$ & $1693\,{:}\,396$ \\
$1697\,{:}\,376$ & $1699\,{:}\,578$ & $1709\,{:}\,1$ & $1721\,{:}\,51$ & $1723\,{:}\,82$ & $1733\,{:}\,248$ & $1741\,{:}\,196$ & $1747\,{:}\,1239$ \\
$1753\,{:}\,432$ & $1759\,{:}\,708$ & $1777\,{:}\,5$ & $1783\,{:}\,106$ & $1787\,{:}\,658$ & $1789\,{:}\,1231$ & $1801\,{:}\,132$ & $1811\,{:}\,259$ \\
$1823\,{:}\,548$ & $1831\,{:}\,992$ & $1847\,{:}\,49$ & $1861\,{:}\,306$ & $1867\,{:}\,71$ & $1871\,{:}\,881$ & $1873\,{:}\,502$ & $1877\,{:}\,250$ \\
$1879\,{:}\,1878$ & $1889\,{:}\,249$ & $1901\,{:}\,76$ & $1907\,{:}\,552$ & $1913\,{:}\,1236$ & $1931\,{:}\,20$ & $1933\,{:}\,1826$ & $1949\,{:}\,1212$ \\
$1951\,{:}\,525$ & $1973\,{:}\,34$ & $1979\,{:}\,1865$ & $1987\,{:}\,586$ & $1993\,{:}\,1$ & $1997\,{:}\,128$ & $1999\,{:}\,6$ & $2003\,{:}\,1268$ \\
$2011\,{:}\,1621$ & $2017\,{:}\,5$ & $2027\,{:}\,1176$ & $2029\,{:}\,80$ & $2039\,{:}\,1150$ & $2053\,{:}\,18$ & $2063\,{:}\,2$ & $2069\,{:}\,997$ \\
$2081\,{:}\,26$ & $2083\,{:}\,19$ & $2087\,{:}\,77$ & $2089\,{:}\,20$ & $2099\,{:}\,1647$ & $2111\,{:}\,1541$ & $2113\,{:}\,1$ & $2129\,{:}\,65$ \\
$2131\,{:}\,116$ & $2137\,{:}\,73$ & $2141\,{:}\,74$ & $2143\,{:}\,58$ & $2153\,{:}\,38$ & $2161\,{:}\,2123$ & $2179\,{:}\,228$ & $2203\,{:}\,73$ \\
$2207\,{:}\,739$ & $2213\,{:}\,133$ & $2221\,{:}\,30$ & $2237\,{:}\,41$ & $2239\,{:}\,9$ & $2243\,{:}\,308$ & $2251\,{:}\,107$ & $2267\,{:}\,118$ \\
$2269\,{:}\,176$ & $2273\,{:}\,237$ & $2281\,{:}\,82$ & $2287\,{:}\,91$ & $2293\,{:}\,112$ & $2297\,{:}\,83$ & $2309\,{:}\,95$ & $2311\,{:}\,24$ \\
$2333\,{:}\,606$ & $2339\,{:}\,588$ & $2341\,{:}\,560$ & $2347\,{:}\,542$ & $2351\,{:}\,1599$ & $2357\,{:}\,506$ & $2371\,{:}\,289$ & $2377\,{:}\,457$ \\
$2381\,{:}\,931$ & $2383\,{:}\,1186$ & $2389\,{:}\,1275$ & $2393\,{:}\,1581$ &  &  &  &  \\
\bottomrule
\end{tabular}
\end{center}

\section{Standalone verification}\label{app:verify}
The script below (requiring only \texttt{sympy}) rebuilds $\N$ from the system of
Appendix~\ref{app:res}, checks the published SHA-256 hashes, confirms that no
prime below $10^{5}$ is a summand, and re-derives the proven example
$\g(m_{43})=89$. It runs to completion with exit code $0$.
\begin{verbatim}
import hashlib
from sympy import isprime, primerange
SHA_N = "853e0013...e729bfc"; SHA_P = "32894e4b...6e367838"
Q, G = 10**5, 100747
RES = [(2,0),(3,2),(5,3),(7,5),(11,8), ...]   # the 356 pairs of Appendix C

def crt(pairs):
    N, M = 0, 1
    for r, a in pairs:
        N += M * (((a - N) * pow(M, -1, r)) % r); M *= r
    return N % M, M

N, M = crt(RES)
if N < 10*Q: N += M
assert hashlib.sha256(str(N).encode()).hexdigest() == SHA_N
assert N % 2 == 0
assert not [q for q in primerange(3, Q)
            if not any((N - q) % r == 0 for (r, a) in RES if r != 2)]
P = N - G
assert isprime(G) and isprime(P)              # G proven; P also ECPP-certified
assert hashlib.sha256(str(P).encode()).hexdigest() == SHA_P
m43 = 1
for p in primerange(2, 44): m43 *= p
assert next(q for q in primerange(2, 1000) if isprime(m43 - q)) == 89
print("ALL CHECKS PASSED")
\end{verbatim}

\noindent The full SHA-256 hashes of the submission artifacts (the decimal
strings for $\N$ and $P$, and the file bytes for the congruence system and the
ECPP certificate) are:
{\small
\begin{verbatim}
N (decimal)        853e0013d3b7627ff5e3a4a814dbc7a346a03ba3df1005fae3dee345ae729bfc
P (decimal)        32894e4bb7c8af3c8a831dfda520fc459f97a7634e696ae8217b23536e367838
residue system     502712ce8342c9998751308190afb4024ebaa5e7edab380bc8ee3225c017d088
ECPP certificate   9ebadbc452ff249fede0af18aee6d3c67c3a101df39a813700235388441edafe
\end{verbatim}}

\end{document}
